\documentclass[12pt, a4paper]{article}

% --- Encoding & Fonts ---
\usepackage[utf8]{inputenc}
\usepackage[T1]{fontenc}
\usepackage{lmodern}

% --- Mathematics ---
\usepackage{amsmath}
\usepackage{amssymb}
\usepackage{amsthm}
\usepackage{mathtools}

% --- Layout ---
\usepackage[margin=1in]{geometry}

% --- Lists ---
\usepackage{enumitem}

% --- Tables ---
\usepackage{booktabs}
\usepackage{array}

% --- Graphics ---
\usepackage{graphicx}

% --- Headers / Footers ---
\usepackage{fancyhdr}
\pagestyle{fancy}
\fancyhf{}
\fancyhead[L]{\leftmark}
\fancyhead[R]{\thepage}
\renewcommand{\headrulewidth}{0.4pt}

% --- Hyperlinks & Cross-references ---
\usepackage[colorlinks=true, linkcolor=blue, urlcolor=blue, citecolor=blue]{hyperref}

% --- Misc ---
\usepackage{xcolor}

% --- Custom commands (safety definitions) ---
\providecommand{\operatorname}[1]{\mathop{\mathrm{#1}}\nolimits}

\title{Learning Module: Janya}
\author{}
\date{\today}

\begin{document}

\maketitle
\thispagestyle{empty}
\newpage

\tableofcontents
\newpage

% ============================================================
% SECTION  Introduction and Foundational Context
% ============================================================

\section{Introduction and Foundational Context}

\subsection{Before We Begin}
A Long ShortTerm Memory (LSTM) network contains three gates---input, forget, and output---that regulate how information is written to, retained in, and read from its internal cell state.  
A Turing machine is an abstract computer that reads and writes symbols on an infinitely long tape, using a head that can move left or right.  
Variablebinding refers to the ability to associate a placeholder (a variable) with a specific value during reasoning, such as linking the name ``X'' to the number 5.  
These ideas together explain why a neural model might still need extra mechanisms to manipulate data flexibly.  
Understanding them lets us see why adding an external, addressable memory can be useful.

\subsection{Motivation}
Think of a chef's kitchen counter: the chef can place ingredients there, glance back at them while cooking, and move them around as needed. The counter acts like a temporary ``working memory'' that the chef constantly reads from and writes to. Standard neural networks lack such an explicit surface to store and retrieve intermediate results, making it hard for them to perform tasks that require stepbystep data manipulation.  

\subsection{Intuition}
Imagine trying to solve a puzzle where you must keep track of several pieces you have already examined while still figuring out where the next piece fits. A regular recurrent network can remember a vague overall sense of the puzzle, but it cannot easily pick out a specific piece, look at it again later, or swap pieces around. By giving the network a dedicated table (the external memory) that it can read from and write to at any time, the network obtains a flexible ``workspace'' much like a human's mental notepad, enabling it to move information around, bind variables to values, and follow explicit logical steps.

\subsection{Formal Definition}
A Recurrent Neural Network (RNN) processes a sequence by updating a hidden state \(\mathbf{h}_t\) at each time step \(t\) based on the current input and the previous hidden state. An RNN is TuringComplete because, with unbounded time and precision, it can simulate the operations of a Turing machine; however, this theoretical capability does not guarantee practical learning of precise read/write procedures.  

An LSTM augments the basic RNN with three gating mechanisms:
- the input gate decides how much new information to add to the cell state,
- the forget gate decides how much of the existing cell state to erase,
- the output gate decides how much of the cell state to expose as the hidden output.  
These gates act as perfect integrators that protect information from vanishing or exploding gradients, but they still operate on a single, fixedsize internal memory vector.  

A Neural Turing Machine (NTM) combines a recurrent controller (often an LSTM) with a memory matrix \(\mathbf{M}\) that contains many memory slots. The controller produces addressing vectors that specify, via differentiable weighting, which slots to read from or write to. Reading returns a weighted sum of the selected slots, while writing can erase and add content to particular slots. This external matrix provides a large, addressable workspace, allowing the system to bind variables by storing their values in distinct locations and to implement explicit logical flow by controlling the sequence of read/write operations.  

Thus, while an RNN can theoretically mimic any algorithm, the NTM supplies a practical mechanism---gated internal state for shortterm retention and a large, addressable external memory for flexible data manipulation---that directly tackles the challenges of logical flow control, explicit working memory, and variablebinding.

\subsection{Key Takeaway}
The Neural Turing Machine equips a recurrent network with an external, addressable memory that supplies the practical workingmemory and variablebinding capabilities that standard RNNs and even LSTMs lack.

\newpage
\subsection{Questions}

% --- RECALL (1) ---
\textbf{Q1.} Define \textit{working memory} as used in cognitive science and give one concrete example.

% --- UNDERSTAND (1) ---
\textbf{Q2.} Explain why standard RNNs, despite being Turingcomplete in theory, are difficult to use for tasks that require explicit data manipulation.

% --- UNDERSTAND (2) ---
\textbf{Q3.} In your own words, describe how the absence of an explicit workspace limits a vanilla RNN when it must keep track of intermediate results during a multistep reasoning process.

% --- APPLY (1) ---
\textbf{Q4.} Consider a simple sequencetosequence copying task where the network must output the input characters in reverse order. Which aspect---logical flow control, external memory, or variable binding---is most likely to cause difficulty for a vanilla LSTM, and why?

% --- APPLY (2) ---
\textbf{Q5.} For a task that requires remembering the position of a particular symbol in a long sequence (e.g., \textit{find the third occurrence of X}), indicate which component (logical flow control, external memory, or variable binding) a vanilla LSTM would struggle with the most.

% --- ANALYSE (1) ---
\textbf{Q6.} Compare the role of LSTM gating mechanisms with the role of the external memory bank in a Neural Turing Machine. Which problems does each mechanism solve?

% --- EVALUATE (1) ---
\textbf{Q7.} Critically assess whether adding a large external memory is the only way to overcome the variablebinding problem in neural networks. Provide arguments supporting and challenging this approach.

\newpage
% ============================================================
% SECTION  Neural Turing Machine Architecture
% ============================================================

\section{Neural Turing Machine Architecture}

\subsection{Before We Begin}
A normalized weighting vector is a list of numbers that are all nonnegative and add up to one.  
Because the numbers sum to one, they can be interpreted as a set of soft probabilities that indicate how much attention is given to each memory location.  
When a head uses such a vector, it can ``look at'' many locations at once, assigning a larger share to some rows and a smaller share to others.  
Even though the head does not pick a single exact cell, the operations remain differentiable because the weighting values change smoothly with the controller's parameters.  
Understanding this ``blurry'' attention prepares us to see how the NTM's controller interacts with its memory bank.

\subsection{Motivation}
Imagine a librarian (the controller) who can both fetch books (read) and place books back on shelves (write).  
The librarian decides which shelf to look at by shining a flashlight that can be broad (blurred) or narrow (sharp).  
This analogy illustrates why a Neural Turing Machine needs a mechanism that can softly focus on multiple memory locations while still being trainable by gradientbased methods.  
The concept solves the problem of giving a neural network a flexible, learnable external memory without sacrificing differentiability.

\subsection{Intuition}
The Neural Turing Machine (NTM) is built from a ``brain'' that decides what to do (the controller) and a notebook of slots that stores information (the memory bank).  
The controller sends out tiny ``beams''---called read heads and write heads---that shine on the notebook.  
Each beam carries a soft focus: instead of pointing at a single slot, it spreads its attention over all slots according to a set of weights that sum to one.  
A read head gathers information by taking a weighted average of the rows it lights up, while a write head updates each row proportionally to its weight, allowing many rows to be affected at once.  
Because the weighting is smooth, tiny changes in the controller's parameters gently shift the focus, keeping the whole system learnable by backpropagation.

\subsection{Formal Definition}
An NTM architecture consists of two principal components:

\begin{itemize}
\item \textbf{Controller}: a neural network that receives external input, produces external output, and generates attention weightings for the memory. The controller may be a recurrent network (e.g., an LSTM) that possesses its own internal state, or a feedforward network that has no internal recurrence.
\item \textbf{Memory bank}: a matrix \(M\) whose rows we call \emph{locations}. Each row stores a vector of fixed length.
\end{itemize}

The controller operates \emph{read heads} and \emph{write heads}.  
For a given head, the controller outputs a normalized weighting vector \(\mathbf{w} = (w_1, w_2, \dots, w_N)\), where \(N\) is the number of rows in \(M\).  Each weight \(w_i\) satisfies \(w_i \ge 0\) and \(\sum_{i=1}^{N} w_i = 1\).  

\paragraph{Read operation.}  
A read head produces a read vector \(\mathbf{r}\) by taking the weighted sum of the rows:
\[
\mathbf{r} = \sum_{i=1}^{N} w_i \, M_{i},
\]
where \(M_{i}\) denotes the \(i\)-th row of the memory matrix.  
Because the weighting is normalized, \(\mathbf{r}\) is a smooth blend of the stored rows; the head can therefore attend sharply to a single location (one weight close to 1) or weakly to many (several moderate weights).

\paragraph{Write operation.}  
A write head supplies a second normalized weighting vector \(\mathbf{w}'\).  For each row \(i\), the content of \(M_{i}\) is updated proportionally to \(w'_i\).  A simple additive write can be expressed as
\[
M_{i} \leftarrow M_{i} + w'_i \, \mathbf{e},
\]
where \(\mathbf{e}\) is the vector emitted by the controller that should be written.  Since the update uses the same smooth weights, the entire write process remains differentiable.

These definitions preempt common misunderstandings: a head does not have to access exactly one cell; a feedforward controller can still influence memory through its heads; and differentiability is achieved precisely because the addressing is soft (blurred), not hard (binary).

\subsection{Worked Example}
\textbf{Problem statement.}  
Consider a memory bank \(M\) with three rows and two columns:
\[
M = \begin{bmatrix}
1 & 2 \\[2pt]
3 & 4 \\[2pt]
5 & 6
\end{bmatrix}.
\]
A feedforward controller produces a readhead weighting vector \(\mathbf{w} = [0.7,\,0.2,\,0.1]\) and a writehead weighting vector \(\mathbf{w}' = [0.1,\,0.6,\,0.3]\).  
The controller also emits a write vector \(\mathbf{e} = [10,\,20]\).  
We will (1) compute the read vector \(\mathbf{r}\) and (2) update the memory using the write operation.

\textbf{Step 1: Compute the read vector.}  
Apply the definition \(\mathbf{r} = \sum_{i=1}^{3} w_i M_i\):
\[
\begin{aligned}
\mathbf{r} &= 0.7 \begin{bmatrix}1 \\ 2\end{bmatrix}
           + 0.2 \begin{bmatrix}3 \\ 4\end{bmatrix}
           + 0.1 \begin{bmatrix}5 \\ 6\end{bmatrix} \\[4pt]
           &= \begin{bmatrix}0.7 \\ 1.4\end{bmatrix}
           + \begin{bmatrix}0.6 \\ 0.8\end{bmatrix}
           + \begin{bmatrix}0.5 \\ 0.6\end{bmatrix} \\[4pt]
           &= \begin{bmatrix}1.8 \\ 2.8\end{bmatrix}.
\end{aligned}
\]

\textbf{Step 2: Apply the write operation.}  
For each row \(i\), update \(M_i \leftarrow M_i + w'_i \mathbf{e}\).

\begin{itemize}
\item Row 1: \(w'_1 = 0.1\)  
\(M_1^{\text{new}} = [1,2] + 0.1\,[10,20] = [1,2] + [1,2] = [2,4]\).

\item Row 2: \(w'_2 = 0.6\)  
\(M_2^{\text{new}} = [3,4] + 0.6\,[10,20] = [3,4] + [6,12] = [9,16]\).

\item Row 3: \(w'_3 = 0.3\)  
\(M_3^{\text{new}} = [5,6] + 0.3\,[10,20] = [5,6] + [3,6] = [8,12]\).
\end{itemize}

Thus the updated memory matrix becomes
\[
M^{\text{new}} = \begin{bmatrix}
2 & 4 \\[2pt]
9 & 16 \\[2pt]
8 & 12
\end{bmatrix}.
\]

\textbf{What this shows us:}  
A read head collects information by averaging rows according to its weighting, while a write head spreads a new vector over all rows in proportion to its own weighting.  Both actions are smooth and differentiable, illustrating the core mechanism of an NTM.

\subsection{Key Takeaway}
An NTM couples a neural controller with a differentiable memory bank; read and write heads use normalized weightings to attend to many locations at once, allowing both recurrent and feedforward controllers to manipulate memory in a trainable way.

\newpage
\subsection*{Questions}

% --- RECALL ---
\textbf{Q1.} List the two main components of a Neural Turing Machine (NTM).

% --- UNDERSTAND ---
\textbf{Q2.} In your own words, explain how a normalized weighting vector allows a read or write head to attend to several memory locations at the same time.

% --- UNDERSTAND ---
\textbf{Q3.} Describe why the use of a normalized weighting vector keeps the read and write operations of an NTM differentiable.

% --- APPLY ---
\textbf{Q4.} A read head produces the weighting vector \(\mathbf{w}= [0.5,\,0.5,\,0]\). The memory matrix is
\[ M = \begin{bmatrix} 1 & 0 \\ 0 & 1 \\ 2 & 2 \end{bmatrix}. \]
Compute the read vector \(\mathbf{r}\) using the definition \(\mathbf{r}=\sum_{i} w_i M_{i}\).

% --- APPLY ---
\textbf{Q5.} Using the same memory matrix \(M\) as in Q4 and the weighting vector \(\mathbf{w}= [0.5,\,0.5,\,0]\), first compute the weighted contribution of each row, then give the final read vector \(\mathbf{r}\). Show your intermediate calculations.

% --- ANALYSE ---
\textbf{Q6.} Contrast the tradeoffs between employing an LSTM controller and a feedforward controller in an NTM. Discuss how each architecture affects the ability to store internal state and how it influences the number of read/write heads required.

% --- EVALUATE ---
\textbf{Q7.} In practical systems, a feedforward controller is said to provide greater transparency but may need additional heads to perform the same tasks as an LSTM controller. Argue whether the benefit of increased transparency outweighs the cost of requiring extra heads, supporting your position with points from the theory.

\newpage
% ============================================================
% SECTION  Reading and Writing Operations
% ============================================================

\subsection*{Notation Used in This Section}
\begin{tabular}{@{}ll@{}}
\toprule
Symbol & Meaning \\ \midrule
    $M_t$ & Memory matrix at time step t (size N\timesM) \\
    $w_t$ & Normalized weighting over memory locations for a head \\
    $r_t$ & Read vector returned by a read head \\
    $e_t$ & Erase vector (elements in (0,1)) \\
    $a_t$ & Add vector injected into memory \\
    $\tilde{M}_t(i)$ & Intermediate memory after erase, before add \\
\bottomrule
\end{tabular}
\vspace{1em}

\section{Reading and Writing Operations}

\subsection{Before We Begin}
A convex combination is a weighted sum where every weight is nonnegative and all weights add up to 1.  
Elementwise multiplication means multiplying two vectors or matrices by matching each of their entries.  
When we talk about a value being in the interval (0, 1) we mean it is greater than 0 and less than 1.  
Understanding these ideas lets us see how a memory cell can be partially emptied and then partially refilled.

\subsection{Motivation}
Mixing paint: you start with a base colour (the memory row), remove a portion of it (erase), then add a new hue (add).  
Just as a painter needs to control how much of each colour is taken away and how much new colour is mixed in, a Neural Turing Machine must be able to read exactly what it needs, erase only the unwanted part of a memory entry, and then add fresh information.  
These operations are designed to be smooth enough for gradientbased learning.

\subsection{Intuition}
The read head looks at all memory locations, decides how important each one is, and then takes a weighted average of those rows -- like blending several shades of paint according to how much of each you want.  
When writing, we first dim the old colour according to an ``erase'' signal, never completely wiping the row, and then spray in a new colour determined by an ``add'' signal.  
Because the dimming (erase) and spraying (add) are independent and can be applied in any order, the final colour depends only on the total amount erased and added, not on the sequence.

\subsection{Formal Definition}
Let \(\mathbf{M}_t\) denote the memory matrix at discrete time step \(t\); it has \(N\) rows (memory locations) and \(M\) columns (the size of each stored vector).  
A read head produces a normalized weighting vector \(\mathbf{w}_t\in\mathbb{R}^N\) whose components satisfy \(\sum_{i=1}^{N} w_t(i)=1\) and \(w_t(i)\ge 0\).  

The read vector \(\mathbf{r}_t\) returned by the head is the convex combination  
\[
\mathbf{r}_t \;\leftarrow\; \sum_{i=1}^{N} w_t(i)\,\mathbf{M}_t(i),
\]
where \(\mathbf{M}_t(i)\) is the \(i\)-th row of \(\mathbf{M}_t\).

Writing proceeds in two stages.  
First, an erase vector \(\mathbf{e}_t\in(0,1)^M\) specifies, elementbyelement, how much of the current content should be removed.  
The intermediate memory after erasing is
\[
\tilde{\mathbf{M}}_t(i) \;\leftarrow\; \mathbf{M}_{t-1}(i)\,\bigl[\,\mathbf{1} - w_t(i)\,\mathbf{e}_t\,\bigr],
\]
where \(\mathbf{1}\) is a length\(M\) vector of ones and the brackets indicate elementwise multiplication.  
Because \(\mathbf{e}_t\) lies in \((0,1)\), each entry of \(\tilde{\mathbf{M}}_t(i)\) is reduced but never driven to a negative value, avoiding a complete deletion of the row.

Second, an add vector \(\mathbf{a}_t\in\mathbb{R}^M\) injects new information:
\[
\mathbf{M}_t(i) \;\leftarrow\; \tilde{\mathbf{M}}_t(i) \;+\; w_t(i)\,\mathbf{a}_t.
\]
The addition is also elementwise; it augments the partiallyerased memory rather than overwriting it.  
If several write heads act on the same memory, the erasethenadd steps commute because both operations are linear in the weighting \(w_t(i)\).

\subsection{Worked Example}
\textbf{Problem statement:}  
Start with a memory matrix having two rows  
\[
\mathbf{M}_0 = \begin{bmatrix}0.8 & 0.2 \\[2pt] 0.1 & 0.9\end{bmatrix}.
\]  
Use weighting \(\mathbf{w} = [0.6,\,0.4]\), erase vector \(\mathbf{e} = [0.3,\,0.7]\), and add vector \(\mathbf{a} = [0.2,\,0.5]\).  
Compute the erased intermediate \(\tilde{\mathbf{M}}_1\), the final memory \(\mathbf{M}_1\), and then the read vector \(\mathbf{r}_1\).

\textbf{Step 1 -- Erase phase (row 1):}  
\[
w_1\mathbf{e} = 0.6\,[0.3,0.7] = [0.18,0.42],
\qquad
\mathbf{1} - w_1\mathbf{e} = [0.82,0.58].
\]  
\[
\tilde{\mathbf{M}}_1(1) = [0.8,0.2]\; \odot\; [0.82,0.58]
          = [0.8\times0.82,\; 0.2\times0.58]
          = [0.656,\;0.116].
\]

\textbf{Step 2 -- Erase phase (row 2):}  
\[
w_2\mathbf{e} = 0.4\,[0.3,0.7] = [0.12,0.28],
\qquad
\mathbf{1} - w_2\mathbf{e} = [0.88,0.72].
\]  
\[
\tilde{\mathbf{M}}_1(2) = [0.1,0.9]\; \odot\; [0.88,0.72]
          = [0.1\times0.88,\; 0.9\times0.72]
          = [0.088,\;0.648].
\]

\textbf{Step 3 -- Add phase (row 1):}  
\[
w_1\mathbf{a} = 0.6\,[0.2,0.5] = [0.12,0.30],
\]  
\[
\mathbf{M}_1(1) = \tilde{\mathbf{M}}_1(1) + w_1\mathbf{a}
               = [0.656,0.116] + [0.12,0.30]
               = [0.776,0.416].
\]

\textbf{Step 4 -- Add phase (row 2):}  
\[
w_2\mathbf{a} = 0.4\,[0.2,0.5] = [0.08,0.20],
\]  
\[
\mathbf{M}_1(2) = \tilde{\mathbf{M}}_1(2) + w_2\mathbf{a}
               = [0.088,0.648] + [0.08,0.20]
               = [0.168,0.848].
\]

\textbf{Step 5 -- Read vector:}  
\[
w_1\mathbf{M}_1(1) = 0.6\,[0.776,0.416] = [0.4656,0.2496],
\]  
\[
w_2\mathbf{M}_1(2) = 0.4\,[0.168,0.848] = [0.0672,0.3392],
\]  
\[
\mathbf{r}_1 = [0.4656,0.2496] + [0.0672,0.3392]
            = [0.5328,0.5888].
\]

\textbf{What this shows us:}  
The erase step only scales down each entry of a memory row according to the weighting and the erase vector, never removing the row entirely.  
The subsequent add step injects new information proportional to the same weighting, and the final read vector is exactly the weighted mixture of the updated rows.

\subsection{Key Takeaway}
Reading extracts a weighted average of memory rows, while writing first dims memory entries with an erase vector (values in (0,1)) and then adds new content; these two phases commute, ensuring stable, differentiable updates.

\newpage
\subsection*{Questions}

{"questions": ["% --- RECALL ---\n\\textbf{Q1.} Write the formula that defines the read vector $\\mathbf{r}_t$ in terms of the weighting vector $\\mathbf{w}_t$ and the rows $\\mathbf{M}_t(i)$ of the memory matrix.", "% --- UNDERSTAND ---\n\\textbf{Q2.} Explain why each component of the erase vector $\\mathbf{e}_t$ must lie strictly in the interval $(0,1)$.", "% --- UNDERSTAND ---\n\\textbf{Q3.} Define a convex combination and describe how the read operation of a Neural Turing Machine forms a convex combination of memory rows.", "% --- APPLY ---\n\\textbf{Q4.} Given a weighting $w = [0.5,\\,0.5]$, an erase vector $e = [0.2,\\,0.4]$, an add vector $a = [0.1,\\,0.3]$, and a single memory row initially $[0.6,\\,0.6]$, compute the updated memory row after the erase and add phases.", "% --- APPLY ---\n\\textbf{Q5.} Using the initial memory matrix $\\mathbf{M}_0 = \\begin{bmatrix}0.8 & 0.2 \\\\ 0.1 & 0.9\\end{bmatrix}$, weighting $\\mathbf{w}= [0.6,\\,0.4]$, erase vector $\\mathbf{e}= [0.3,\\,0.7]$ and add vector $\\mathbf{a}= [0.2,\\,0.5]$, compute (a) the intermediate erased matrix $\\tilde{\\mathbf{M}}_1$, (b) the final memory matrix $\\mathbf{M}_1$, and (c) the read vector $\\mathbf{r}_1$.", "% --- ANALYSE ---\n\\textbf{Q6.} If two write heads use the same weighting vector but different erase vectors, discuss how their combined effect determines the final content of each memory row.", "% --- ANALYSE ---\n\\textbf{Q7.} Suppose two write heads apply different weighting vectors while sharing the same erase and add vectors. Explain whether the order in which the heads are applied changes the final memory matrix, and justify your answer using the commutativity property.", "% --- EVALUATE ---\n\\textbf{Q8.} Assess the advantages of separating the erase and add phases in a write operation compared to performing a single overwrite of a memory row."}

\newpage
% ============================================================
% SECTION  Addressing Mechanisms
% ============================================================

\subsection*{Notation Used in This Section}
\begin{tabular}{@{}ll@{}}
\toprule
Symbol & Meaning \\ \midrule
    $w_t$ & Final weighting over memory locations at time t \\
    $w_t^c$ & Contentbased weighting \\
    $k_t$ & Key vector emitted by controller \\
    $\beta_t$ & Key strength (sharpness) scalar \\
    $K[\cdot,\cdot]$ & Cosine similarity between two vectors \\
    $g_t$ & Interpolation gate scalar \\
    $w_t^g$ & Gated weighting (interpolated) \\
    $s_t$ & Shift weighting defining allowed integer shifts \\
    $\tilde{w}_t$ & Weighting after circular convolution (shifted) \\
    $\gamma_t$ & Sharpening scalar (\geq1) \\
\bottomrule
\end{tabular}
\vspace{1em}

\section{Addressing Mechanisms}

\subsection{Before We Begin}
Cosine similarity measures how alike two vectors are by looking at the angle between them.  
A softmax converts a list of numbers into a probabilitylike distribution by exponentiating each number and then normalising by the sum of all exponentials.  
Basic circular convolution shifts a vector by mixing its entries according to a small shift pattern; it is essentially a ``wraparound'' shift.  
A scalar gate is a single number between 0 and 1 that decides how much of one quantity to keep versus another.  
Understanding these tools will let us see how the NTM builds its final weighting over memory locations.

\subsection{Motivation}
Finding a book in a library: you first look for the cover color that matches (contentbased), then you might move one shelf left or right (location shift), and finally you focus more sharply on the exact spot.  
The NTM needs a systematic way to decide which memory cell to read or write, combining what looks most relevant with where it is in memory.  
By blending similaritydriven focus with controlled shifts, the model can both retrieve specific items and walk through memory sequentially.

\subsection{Intuition}
At each time step the NTM asks two questions: ``What does the data I'm looking for look like?'' and ``Where should I look relative to where I was before?''  
The first answer produces a soft focus on all memory slots that resemble a query vector; the second answer moves that focus a little left or right (or leaves it unchanged).  
After the shift the distribution may become blurry, so the system sharpens it again, ending with a clear set of weights that indicate which slot(s) to attend to.

\subsection{Formal Definition}
Let \(w_t\) denote the final weighting over the \(N\) memory locations at time step \(t\).

\paragraph{Contentbased weighting}
The controller emits a key vector \(k_t\).  
For each memory vector \(M_t(i)\) we compute the cosine similarity
\[
K[k_t,\,M_t(i)] \;=\; \frac{k_t\!\cdot\! M_t(i)}{\|k_t\|\,\|M_t(i)\|}.
\]
A keystrength scalar \(\beta_t\) sharpens the similarities, and a softmax yields the content weighting
\[
w_t^{c}(i)\;=\;
\frac{\exp\!\big(\beta_t\,K[k_t,\,M_t(i)]\big)}
     {\displaystyle\sum_{j=1}^{N}\exp\!\big(\beta_t\,K[k_t,\,M_t(j)]\big)} .
\]
The larger \(\beta_t\) is, the more the distribution concentrates on the most similar locations; a small \(\beta_t\) spreads attention more evenly.

\paragraph{Interpolation}
An interpolation gate \(g_t\in[0,1]\) mixes the new content weighting with the previous weighting \(w_{t-1}\):
\[
w_t^{g}\;=\; g_t\,w_t^{c}\;+\;(1-g_t)\,w_{t-1}.
\]
When \(g_t\) is close to 1 the model trusts the fresh content focus, while a value near 0 lets the old focus dominate; even intermediate values retain contributions from both sources.

\paragraph{Locationbased addressing (circular convolution)}
A shift weighting \(s_t\) is a normalized distribution over allowed integer shifts (for example \([-1,0,1]\)).  
The gated weighting is shifted by circular convolution:
\[
\tilde{w}_t(i)\;=\;\sum_{j=0}^{N-1} w_t^{g}(j)\;s_t(i-j),
\]
where indices wrap around, so a shift that would move past the last location reappears at the first.

\paragraph{Sharpening}
Because convolution can smear the distribution, a sharpening scalar \(\gamma_t\ge 1\) reconcentrates the mass:
\[
w_t(i)\;=\;
\frac{\tilde{w}_t(i)^{\gamma_t}}
     {\displaystyle\sum_{j=1}^{N}\tilde{w}_t(j)^{\gamma_t}} .
\]
Raising each entry to a power \(\gamma_t>1\) makes large values relatively larger and small values relatively smaller, producing a sharper, more peaked final weighting.

These steps can be applied together, so contentbased and locationbased addressing are not mutually exclusive; the interpolation gate never completely discards the previous weighting unless \(g_t=0\); and sharpening makes the distribution more peaked, not more uniform.

\subsection{Worked Example}
\textbf{Problem statement}\\
The memory contains three rows:
\[
M(1)=[0.5,0.1],\quad
M(2)=[0.6,0.2],\quad
M(3)=[0.2,0.8].
\]
The controller emits the key vector \(k=[0.6,0.2]\).  
We will compute the final weighting \(w_t\) using \(\beta_t=5\), previous weighting \(w_{t-1}=[0.2,0.5,0.3]\), interpolation gate \(g_t=0.6\), shift weighting \(s_t=[0,1,0]\) (a rightshift by one), and sharpening scalar \(\gamma_t=2\).

\paragraph{Step 1: Cosine similarities}
\[
\begin{aligned}
\|k\| &= \sqrt{0.6^2+0.2^2}=0.6325,\\
\|M(1)\| &= \sqrt{0.5^2+0.1^2}=0.5099, &
K[k,M(1)] &= \dfrac{0.6\cdot0.5+0.2\cdot0.1}{0.6325\cdot0.5099}\approx0.994,\\[4pt]
\|M(2)\| &= \sqrt{0.6^2+0.2^2}=0.6325, &
K[k,M(2)] &= \dfrac{0.6\cdot0.6+0.2\cdot0.2}{0.6325^2}=1.000,\\[4pt]
\|M(3)\| &= \sqrt{0.2^2+0.8^2}=0.8246, &
K[k,M(3)] &= \dfrac{0.6\cdot0.2+0.2\cdot0.8}{0.6325\cdot0.8246}\approx0.537.
\end{aligned}
\]

\paragraph{Step 2: Apply key strength \(\beta_t=5\) and softmax}
Unnormalised exponentials:
\[
\begin{aligned}
e^{5\cdot0.994}&\approx e^{4.97}\approx144.6,\\
e^{5\cdot1.000}&=e^{5}\approx148.4,\\
e^{5\cdot0.537}&\approx e^{2.685}\approx14.66.
\end{aligned}
\]
Normalization constant \(Z=144.6+148.4+14.66\approx307.66\).

\[
w_t^{c}\approx\bigl[0.470,\;0.483,\;0.048\bigr].
\]

\paragraph{Step 3: Interpolation with previous weighting}
\[
w_t^{g}=0.6\,w_t^{c}+0.4\,w_{t-1}
\]
\[
\begin{aligned}
w_t^{g}(1)&=0.6(0.470)+0.4(0.2)=0.362,\\
w_t^{g}(2)&=0.6(0.483)+0.4(0.5)=0.4898,\\
w_t^{g}(3)&=0.6(0.048)+0.4(0.3)=0.1488.
\end{aligned}
\]
(These three numbers sum to approximately 1.)

\paragraph{Step 4: Circular convolution with shift \(s_t=[0,1,0]\)}
Because \(s_t\) places all mass on a shift of +1, the operation simply rotates the vector right by one position:
\[
\tilde{w}_t=\bigl[\,w_t^{g}(3),\;w_t^{g}(1),\;w_t^{g}(2)\,\bigr]
          \approx\bigl[0.1488,\;0.362,\;0.4898\bigr].
\]

\paragraph{Step 5: Sharpening with \(\gamma_t=2\)}
\[
\begin{aligned}
\tilde{w}_t(1)^{2}&=0.1488^{2}=0.02215,\\
\tilde{w}_t(2)^{2}&=0.362^{2}=0.131,\\
\tilde{w}_t(3)^{2}&=0.4898^{2}=0.240.
\end{aligned}
\]
Normalization constant \(Z'=0.02215+0.131+0.240\approx0.393\).

\[
w_t\approx\bigl[0.056,\;0.333,\;0.611\bigr].
\]

\textbf{What this shows us:} the NTM first focuses on the memory slot most similar to the key, then moves that focus one step to the right, and finally sharpens the result so that the attention concentrates strongly on a single location.

\subsection{Key Takeaway}
The NTM builds its final memory weighting by (1) computing a similaritybased soft focus, (2) interpolating with the previous focus, (3) shifting that focus with a circular convolution, and (4) sharpening the outcome to obtain a clear, peaked distribution.

\newpage
\subsection*{Questions}

% --- RECALL ---
\\textbf{Q1.} Write the equation that defines the contentbased weighting $w_t^{c}(i)$ in terms of the key strength $\beta_t$, cosine similarity $K[k_t,M_t(i)]$, and the softmax normalisation.

% --- UNDERSTAND ---
\\textbf{Q2.} Explain the role of the keystrength scalar $\beta_t$ in the softmax computation over cosine similarities. How does changing $\beta_t$ affect the resulting contentbased weighting?

% --- UNDERSTAND ---
\\textbf{Q3.} In your own words, describe why the NTM combines contentbased addressing with an interpolation gate before applying a locationbased shift. What advantage does this give the model?

% --- APPLY ---
\\textbf{Q4.} The cosine similarities between a key vector and three memory locations are $[0.8,\,0.4,\,0.2]$ and the keystrength is $\beta_t = 3$. Compute the unnormalised exponentials $e^{\beta_t K}$ for each location.

% --- APPLY ---
\\textbf{Q5.} Using the same data as in the previous question, compute the normalised contentbased weighting $w_t^{c}(i)$ by applying the softmax (divide each unnormalised exponential by the sum of all three). Give the resulting three values rounded to three decimal places.

% --- ANALYSE ---
\\textbf{Q6.} If the interpolation gate $g_t$ is set to $0$, what becomes of the new gated weighting $w_t^{g}$? Discuss how this choice would influence the NTM's ability to learn a new memory location.

% --- EVALUATE ---
\\textbf{Q7.} The sharpening step raises each entry of $\tilde{w}_t$ to the power $\gamma_t\ge 1$ before renormalising. Critique the use of this sharpening after the circular convolution. Could the NTM operate effectively without the sharpening step? Provide reasons based on the behaviour of convolution and the need for a peaked distribution.

\newpage
% ============================================================
% SECTION  Experiments and Algorithmic Tasks
% ============================================================

\section{Experiments and Algorithmic Tasks}

\subsection{Before We Begin}
An episodic task is a learning problem in which the internal state of the model is reset at the beginning of each input sequence, so that each sequence is processed independently.  During training the model sees sequences up to a certain maximum length, while during testing it may be asked to handle longer sequences that were never seen before.  To ``generalise'' means that the model can correctly handle these longer test sequences even though it was only trained on shorter examples.  Understanding these ideas lets us see why the results reported for the NTM are significant.

\section{Motivation}
Imagine sorting a pile of mail by urgency: you first drop each letter into a slot labelled with its priority (low, medium, high). Later you walk along the slots in order and collect the letters, producing a sorted batch.  Researchers created a series of small algorithmic problems---copying a sequence, repeating a copy, associative recall, dynamic Ngrams, and priority sort---to ask whether a neural model can learn the underlying procedures rather than merely memorising examples.  The goal is to reveal whether the Neural Turing Machine can acquire looplike and tablelike behaviours that let it work on inputs far larger than those used for training.

\section{Intuition}
The NTM augments a neural controller with an external memory that can be written to and read from at each time step.  In the copy task, the controller stores each incoming symbol in a fresh memory location, then later steps through those locations in order to reproduce the sequence.  Because the same writeoncereadonce pattern can be repeated arbitrarily many times, the NTM behaves like a simple loop that does not depend on the length it was trained on.  Similar ideas underlie the other tasks: a nested loop can repeat a copy, a contentbased lookup plus a fixed offset can retrieve the next item, a rewritable table can count Ngram frequencies, and writing according to a numeric priority yields a sorted readout.

\section{Formal Definition}
Consider a set of five benchmark tasks:

\begin{enumerate}
\item \textbf{Copy}:  For an input sequence $[x_{1},x_{2},\dots ,x_{T}]$ the model must output the identical sequence.  The NTM writes each $x_{t}$ to a distinct memory slot $M_{t}$ and later reads from $M_{1},M_{2},\dots ,M_{T}$ in that order.
\item \textbf{Repeat Copy}:  The model receives an input sequence together with an integer $R$ indicating how many times the sequence should be reproduced.  The NTM executes an outer loop $R$ times, each time iterating through the previously stored memory slots.
\item \textbf{Associative Recall}:  A list of items $[a_{1},a_{2},\dots ,a_{N}]$ is presented, followed by a query item $q$.  The required output is the item that immediately follows $q$ in the list.  The NTM stores a compressed representation of each $a_{i}$, performs a contentbased lookup to locate the memory slot containing $q$, then reads the slot that is one position forward.
\item \textbf{Dynamic NGrams}:  At each time step the model observes a symbol from a stream generated by an unknown $6$gram distribution and must predict the next symbol.  The NTM treats its memory as a table that records counts of observed $6$grams; this table can be overwritten as new data arrive, thereby implementing a Bayesianlike estimator.
\item \textbf{Priority Sort}:  Each input vector arrives with an associated scalar priority $p\in[0,1]$.  The NTM writes the vector into a memory address proportional to $p$ (e.g., low $p$ near the beginning, high $p$ near the end).  After all inputs are stored, the read head scans the memory from low to high addresses, producing the vectors sorted by priority.
\end{enumerate}

In the copy and repeatcopy tasks the NTM's behavior reproduces an \emph{arrayiteration loop} and a \emph{nested loop}, respectively.  In associative recall, the combination of \emph{contentbased addressing} (finding the slot that matches the query) and a fixed \emph{location offset} (moving one slot forward) yields the correct subsequent item.  For dynamic Ngrams, the external memory functions as a \emph{rewritable frequency table}, allowing the model to adapt instantly to a new distribution.  In priority sort, the strategy of writing according to the scalar priority and then reading sequentially implements a genuine sorting procedure rather than a fixed mapping.

\section{Worked Example}
\textbf{Problem statement:}  Perform the copy task on the short input sequence $[a,b,c]$ and show that the same steps work for the longer, unseen sequence $[d,e,f,g,h]$.

\textbf{Step 1 -- Initialization.}  
At time $t=0$ the memory is empty and the read and write heads are ready to address the first location.

\textbf{Step 2 -- Write phase for $[a,b,c]$.}  
\begin{itemize}
\item $t=1$: the controller receives symbol $a$, writes $a$ to memory slot $M_{1}$, and moves the write head to the next slot.
\item $t=2$: the controller receives symbol $b$, writes $b$ to $M_{2}$, and moves the write head forward.
\item $t=3$: the controller receives symbol $c$, writes $c$ to $M_{3}$, and moves the write head forward.
\end{itemize}

\textbf{Step 3 -- Read phase for $[a,b,c]$.}  
\begin{itemize}
\item $t=4$: the read head points to $M_{1}$, reads $a$, and outputs $a$.
\item $t=5$: the read head moves to $M_{2}$, reads $b$, and outputs $b$.
\item $t=6$: the read head moves to $M_{3}$, reads $c$, and outputs $c$.
\end{itemize}

\textbf{Step 4 -- Write phase for the longer sequence $[d,e,f,g,h]$.}  
The same rule (``write the current input symbol to the next free memory slot'') is applied:
\begin{itemize}
\item $t=1$: write $d$ to $M_{1}$.
\item $t=2$: write $e$ to $M_{2}$.
\item $t=3$: write $f$ to $M_{3}$.
\item $t=4$: write $g$ to $M_{4}$.
\item $t=5$: write $h$ to $M_{5}$.
\end{itemize}

\textbf{Step 5 -- Read phase for $[d,e,f,g,h]$.}  
\begin{itemize}
\item $t=6$: read $M_{1}$ \rightarrow output $d$.
\item $t=7$: read $M_{2}$ \rightarrow output $e$.
\item $t=8$: read $M_{3}$ \rightarrow output $f$.
\item $t=9$: read $M_{4}$ \rightarrow output $g$.
\item $t=10$: read $M_{5}$ \rightarrow output $h$.
\end{itemize}

\textbf{Final answer:}  The NTM reproduces $[a,b,c]$ correctly and, without any change to its internal algorithm, also reproduces the unseen longer sequence $[d,e,f,g,h]$.

\textbf{What this shows us:}  The copy task does not rely on memorising a fixedlength pattern; instead the NTM implements a loop that writes each input to a fresh memory cell and later reads the cells sequentially.  Because the loop rule is independent of the sequence length, the model generalises to arbitrarily long inputs.

\subsection{Key Takeaway}
The NTM learns algorithmic looping and memorymanipulation strategies that let it handle sequences far longer than those seen during training, whereas a standard LSTM degrades because it lacks such explicit external memory operations.

\newpage
\subsection*{Questions}

% --- RECALL (1) ---
\\textbf{Q1.} List the five benchmark algorithmic tasks that were used to evaluate the Neural Turing Machine (NTM) in the experiments described.

% --- UNDERSTAND (1) ---
\\textbf{Q2.} In your own words, explain why the NTM is able to generalise almost perfectly to longer sequences in the copy task, whereas a standard LSTM's performance declines rapidly.

% --- UNDERSTAND (2) ---
\\textbf{Q3.} Which specific aspect of the NTM's architecture (as described in the theory) enables it to implement a lengthindependent loop for copying, and why is this advantage not present in a plain LSTM?

% --- APPLY (1) ---
\\textbf{Q4.} Suppose the input sequence is $[x_{1},x_{2}]$. Describe stepbystep what the NTM does during the write phase (including which memory slots are used) and then during the read phase to reproduce the same twosymbol sequence.

% --- APPLY (2) ---
\\textbf{Q5.} Using the timestep notation from the worked example, fill in the actions (write or read, symbol, memory slot) that the NTM would perform for the input sequence $[p,q]$ from $t=1$ through $t=6$.

% --- ANALYSE (1) ---
\\textbf{Q6.} Compare how the NTM and a nave LSTM would approach the repeatcopy task. Discuss the differences in memory usage, looping mechanism, and expected ability to handle a larger number of repetitions than seen during training.

% --- EVALUATE (1) ---
\\textbf{Q7.} Critically assess whether the prioritysort experiment provides evidence that the NTM has learned a genuine sorting algorithm or merely a deterministic mapping from priority to memory address. Use the described strategy of writing vectors to memory locations proportional to their scalar priority to support your argument.

\end{document}
